% This is samplepaper.tex, a sample chapter demonstrating the
% LLNCS macro package for Springer Computer Science proceedings;
% Version 2.20 of 2017/10/04
%
\documentclass[runningheads]{llncs}
%
\usepackage{lipsum}  
\usepackage{multirow}
% \usepackage{bbding}
\usepackage{graphicx}
\usepackage{marvosym}
\usepackage{amsfonts,amssymb}
\usepackage[export]{adjustbox}
\usepackage{pdfpages}
\newcommand{\tabincell}[2]{\begin{tabular}{@{}#1@{}}#2\end{tabular}}
% Used for displaying a sample figure. If possible, figure files should
% be included in EPS format.
%
% If you use the hyperref package, please uncomment the following line
% to display URLs in blue roman font according to Springer's eBook style:
% \renewcommand\UrlFont{\color{blue}\rmfamily}

\begin{document}
%
% \title{Conditional Flow Variational Autoencoder for Meaningful Virtual Population Modeling}
\title{A Conditional Flow Variational Autoencoder for Controllable Synthesis of Virtual Populations of Anatomy}
%
%\titlerunning{Abbreviated paper title}
% If the paper title is too long for the running head, you can set
% an abbreviated paper title here
%
\author{Haoran Dou\inst{1} \and
Nishant Ravikumar\inst{1,2}\thanks{Nishant Ravikumar and Alejandro F. Frangi are joint last author} \and
Alejandro F. Frangi\inst{1,2,3,4\star}\textsuperscript{(\Letter)}}
%
\authorrunning{Haoran Dou et al.}
% First names are abbreviated in the running head.
% If there are more than two authors, 'et al.' is used.
%
\institute{Centre for Computational Imaging and Simulation Technologies in Biomedicine (CISTIB), University of Leeds, Leeds, UK \and
NIHR Leeds Biomedical Research Centre (BRC), Leeds, UK \and
Alan Turing Institute, London, UK \and
Medical Imaging Research Center (MIRC), Electrical Engineering and Cardiovascular Sciences Departments, KU Leuven, Leuven, Belgium
% Springer Heidelberg, Tiergartenstr. 17, 69121 Heidelberg, Germany
\\
\email{\{a.frangi, n.ravikumar\}@leeds.ac.uk}}
%
\maketitle              % typeset the header of the contribution
%
\begin{abstract}
Generating virtual populations (VPs) of anatomy is essential for conducting in-silico trials of medical devices. Typically, the generated VP should capture sufficient variability while remaining plausible, and should reflect specific characteristics and patient demographics observed in real populations. It is desirable in several applications to synthesize VPs in a \textit{controlled} manner, where relevant covariates are used to conditionally synthesise virtual populations that fit specific target patient populations/characteristics. %Recent techniques on \textit{controllable} VP synthesis for anatomical models are usually based on conditional variational auto-encoders (cVAEs) with an assumed Gaussian prior having diagonal unit variance. However, such a unimodal prior fails to model multi-modal latent distributions and hence, fails to capture the observed anatomical variability sufficiently. Therefore, 
We propose to equip a conditional variational autoencoder (cVAE) with normalizing flows to boost the flexibility and complexity of the approximate posterior learned, leading to enhanced flexibility for controllable synthesis of VPs of anatomical structures. We demonstrate the performance of our conditional-flow VAE using a dataset of cardiac left ventricles acquired from 2360 patients, with associated demographic information and clinical measurements (used as covariates/conditioning information). %Meanwhile, we manipulated the data to validate the model's ability to capture the correlations between demographics and anatomies qualitatively. 
The obtained results indicate the superiority of the proposed method for conditional synthesis of virtual populations of cardiac left ventricles relative to a cVAE. Conditional synthesis performance was assessed in terms of generalisation and specificity errors, and in terms of the ability to preserve clinical relevant biomarkers in the synthesised VPs, I.e. left ventricular blood pool and myocardial volume, relative to the observed real population.%, in generating a meaningful and diverse VP compared with vanilla cVAE.
\keywords{Virtual Population  \and Generative Model \and Normalizing Flow.}
\end{abstract}
%
%
%
\section{Introduction}
\textit{In-silico} trials (ISTs) use computational modelling and simulation techniques with virtual twin or patient models of anatomy and physiology to evaluate the safety and efficacy of medical devices virtually~\cite{viceconti2021silico}. Virtual patient populations (VPs), distinct from virtual twin populations, comprise plausible instances of anatomy and physiology that do not represent any specific real patient's data (as in the case of the latter, viz. virtual twins). In other words, VPs comprise synthetic data, that help expand/enrich the diversity of anatomical and physiological characteristics that can be investigated within an IST for a given medical device. %This in turn can help generate more comprehensive evidence and distill deeper insights as to the performance of a medical device, especially in operating conditions (i.e. anatomical/physiological configurations) that are under-represented in real patient populations.
A key aspect of patient recruitment in real clinical trials used to assess device performance and generate regulatory evidence for the approval of devices, is the clear definition of inclusion and exclusion criteria for the trial. These criteria define the target patient population deemed appropriate/safe to assess the performance of the device of interest. Correspondingly, it is therefore desirable to enable \textit{controlled} synthesis of VPs that may be used for ISTs of devices, in a manner that emulates imposition of trial inclusion and exclusion criteria. %Achieving controllable synthesis of meaningful and diverse VPs of anatomy is a fundamental yet challenging step to deliver valid and comprehensive scientific evidence in ISTs. 
%Typically, the inclusion and exclusion criteria, such as age, gender, ethnicity, etc., will be defined in the trial protocol when recruiting patients to a real clinical trial, thus ensuring the target population satisfies the research purpose. This proposes a high requirement for controllable VPs synthesis to ensure the VPs could reflect specific patient characteristics and demographics. 

% in digital twins or virtual patient populations (VPs) to provide evidence for evaluating the effectiveness and safety of the medical device virtually~\cite{viceconti2021silico}. 

Virtual populations can be considered to be parametric representations of the anatomy sampled from a generative model. Traditional statistical shape models (SSMs), based on methods such as principal component analysis (PCA), have been widely explored in the last decade~\cite{ravikumar2018group,gooya2015bayesian,frangi2002automatic}. Recent studies focus on deep-learning-based generative models due to their automatic and powerful hierarchical feature extraction~\cite{bonazzola2021image,dou2022generative}. For instance, Bonazzola~\textit{et al.}~\cite{bonazzola2021image} utilized a graph convolutional Varational Auto-encoder (gcVAE) to learn the latent representations of the 3D left ventricular meshes and used the learned representations as surrogates for cardiac phenotypes in genome-wide association studies. Dou~\textit{et al.}~\cite{dou2022generative} proposed learning the shape representations of multiple cardiovascular anatomies via gcVAEs independently and then assembling them into complete whole-heart anatomies termed virtual heart chimaeras. Other studies have investigated conditional generative models for synthesising VPs of anatomies. For example, Beetz~\textit{et al.}~\cite{beetz2022generating} employed a conditional VAE (cVAE), conditioned on gender and cardiac phase, to enable synthesis of VPs of biventricular anatomies. In subsequent work~\cite{li2023deep,beetz2022multi}, the authors extended their method to a multi-domain VAE for modelling biventricular anatomies at multiple times points (across the cardiac cycle), by utilising patient-specific electrocardiogram (ECG) signals as additional (on top of patient demographic data and standard clinical measurements) conditioning information to guide the synthesis. %The electrocardiogram (ECG) was also leveraged as an additional condition. 
%Another study used a recurrent neural network to capture spatio-temporal patterns of biventricular shape and motion observed across a population within the latent representations learned by a cVAE, where, multiple clinical measurements and patient demographic information were used as covariates for conditional, sequential synthesis of biventricular shape and motion~\cite{qiao2023cheart}. 
All aforementioned methods model the latent space in the VAEs/cVAEs as a multivariate Gaussian distribution with a diagonal covariance matrix. This limits the flexibility afforded to the cVAE, as the Gaussian distribution, being unimodal, is a poor approximation to multi-modal latent posterior distributions. This in turn limits the overall variability in anatomical shape that can be captured by standard VAEs and cVAEs. %Hence, enriching the latent posterior distribution to boost the diversity and plausibility of generative VPs remains challenging.

In this study, we address the limitations of existing state-of-the-art conditional generative models used synthesise VPs of anatomical structures. In particular, we propose a method to relax the constraint on modelling the latent distribution as a unimodal multi-variate Gaussian, to boost the flexibility of the generative model and enable conditional synthesis of diverse and plausible VPs generation. Recent advances in normalizing flows~\cite{rezende2015variational,tomczak2016improving} introduce a new solution for this limitation by leveraging a series of invertible parameterized functions to transform the unimodal distribution to a multi-modal one. Motivated by this technique, we propose the first conditional flow VAE (parameterised as a graph-convolutional network) for the task of \textit{controllable} synthesis of VPs of anatomy. The contributions are as follows: (i) we introduce normalizing flows to learn a multi-modal latent posterior distribution by transforming the latent variables from a simple unimodal distribution. This helps the generative model capture greater anatomical variability from the observed real population, leading to the synthesis of more diverse VPs; (ii) we condition the flow-based VAE on patient demographic data and clinical measurements. This enables conditional synthesis of plausible VPs (given relevant covariates/conditioning information as inputs), which reflect the observed correlations between non-imaging patient information and anatomical characteristics in the real population.

\section{Methodology}
In this study, we propose a cVAE model equipped with normalizing flows for the controllable synthesis of VPs of cardiovascular anatomy. A schematic of the proposed conditional flow VAE network's architecture is shown in Figure 1. We employ normalizing flows in the latent space of the cVAE to transform the initial Gaussian posterior to a complex multimodal distribution. 

\begin{figure*}[!htbp]
	\centering
	\includegraphics[width=0.9\linewidth]{Framework-cropped.pdf}	
	\caption{Schematic illustration of our proposed conditional flow VAE}
	\label{fig:framework}
\end{figure*}

% \subsection{Conditional Variational Autoencoder}
% \label{sec:cvae}
\textbf{Conditional Variational Autoencoder:} A VAE is a probabilistic generative model/network~\cite{kingma2013auto} that comprises an encoder and a decoder network branch. The encoder learns a mapping from the input data to a low-dimensional latent space that abstracts the semantic representations from the observations, and the decoder reconstructs the original data from the low-dimensional latent representation. The latent space from which the observed data is generated, is given by approximating the posterior distribution of the latent variables using variational inference. The VAE network is trained by maximizing the evidence lower bound (ELBO), which is a summation of the expected log-likelihood of the data and the Kullback-Leibler divergence between the approximate posterior and some assumed prior distribution over the latent variables (typically a multi-variate Gaussian distribution). Despite its effectiveness at capturing some of the observed variability in the training population (e.g. of anatomical shapes or images), VAEs do not provide any control over the generation process and hence cannot guarantee that the generated population anatomical shapes are representative of target patient populations with specific inclusion/exclusion criteria. Controllable synthesis of anatomical VPs is essential for constructing meaningful cohorts for use in ISTs. Conditional VAE~\cite{sohn2015learning} is a VAE-variant that uses additional covariates/conditioning information in addition to the input data (e.g. anatomical shapes) to learn a conditional latent posterior distribution (conditioned on the covariates), enabling controllable synthesis of VPs during inference (given relevant covariates/conditioning information as inputs).

Our conditional flow VAE (cVAE-NF) is a graph-convolutional network which takes as input a triangular surface mesh representation of an anatomical structure of interest, i.e., the Left Ventricle (LV) in this study, and its associated covariates/conditioning variables, i.e., the patient demographic data and clnical measurements, such as gender, age, weight, blood cholesterol, etc., and outputs the reconstructed surface mesh. Each mesh is represented by a list of 3D spatial coordinates of its vertices and an adjacency matrix defining vertex connectivity (i.e. edges of mesh triangles). The encoder and decoder contain five residual graph-convolutional blocks, respectively. Each block comprises two Chebyshev graph convolutions, each of which is followed by batch normalisation, and ELU activation. A residual connection is added between the input and the output of each graph-convolutional block. Hierarchical mesh down/up-sampling operations proposed in CoMA~\cite{ranjan2018generating} are adopted after each block to capture the global and local shape context. The VAE model is conditioned on the covariates by scaling the hidden representations in the encoder similar to adaptive instance normalization~\cite{huang2017arbitrary} given the covariates as input to generate the scaling factor, and by concatenating the covariates with the latent variables before decoding.

% \subsection{Flexible Posterior using Normalizing Flow}
% \label{sec:nf}
\begin{figure*}[!htbp]
	\centering
	\includegraphics[width=0.8\linewidth]{Flow-cropped.pdf}	
	\caption{Effect of normalizing flow on Gaussian distribution. Step 0 is the initial two-dimensional Gaussian distribution, and step 1-5 represents the distribution of latent variables transformed by the normalizing flow layers (i.e., planar flow).}
	\label{fig:nf}
\end{figure*}

\textbf{Flexible Posterior using Normalizing Flow:} Vanilla cVAEs model the approximate posterior distribution using Gaussian distributions with a diagonal covariance matrix. However, such a unimodal distribution is a poor approximation of the complex true latent posterior distribution in most real-world applications (e.g. for shapes of the LV observed across a population), limiting the anatomical variability captured by the model. In this study, we introduce normalizing flows to construct a flexible multi-modal latent posterior distribution by applying a series of differentiable, invertible/diffeomorphic transformations iteratively to the initial simple unimodal latent distribution. As shown in Fig~\ref{fig:nf}, a two-dimensional Gaussian distribution can be transformed into a multi-modal distribution by applying several normalizing flow steps to the former.

Consider an invertible and smooth mapping function $f: \mathbb{R}^{d} \rightarrow \mathbb{R}^{d}$ with inverse $f^{-1}=g$, and a random variable $\textbf{z}$ with distribution $q(\textbf{z})$. The transformed variable $\textbf{z}^{\prime}=f(\textbf{z})$ follows a distribution given by: 
\begin{equation}
    q( \textbf{z}^{\prime}) = q(\textbf{z}) \left\vert {\rm det} \frac{\partial f}{\partial \textbf{z}}  \right\vert^{-1}
\end{equation}
where the ${\rm det} \frac{\partial f}{\partial \textbf{z}}$ is the Jacobian determinant of $f$. Therefore, we can obtain a complex multi-modal density by composing multiple invertible mappings to transform the initial, simple and tractable density sequentially, as follows,
\begin{equation}
    \textbf{z}_{i} = f_{i}\circ \dots \circ f_{2} \circ f_{1}(\textbf{z}_{0})
\end{equation}
\begin{equation}
\label{equ:lnqz}
    \ln q_{i}(\textbf{z}_{i}) = \ln q_{0}(\textbf{z}_{0}) - \sum^{i} \ln \left\vert {\rm det} \frac{\partial f_{i}}{\partial \textbf{z}_{i-1}} \right\vert
\end{equation}

The specific mathematical formulation of the normalizing flow function is important and must be chosen with care to allow for efficient gradient computation during training, scalable inference and efficiency in computing the determinant of the Jacobian. In this study, we leverage the planar flow in~\cite{rezende2015variational} as a basic unit of our latent normalizing flow net. Specifically, each transformation unit is given by,
\begin{equation}
    f(\textbf{z}) = \textbf{z} + \textbf{u} h(w^{\top}\textbf{z} + b)
\end{equation}
where $\textbf{w} \in \mathbb{R}^d$, $\textbf{u} \in \mathbb{R}^d$ and $b\in \mathbb{R}$ are learnable parameters; $h(\cdot)$ is a smooth element-wise non-linear function with derivative $h^{\prime}(\cdot)$ (we use $\tanh$ in our study) and $\mathbf{z}$ denotes the latent variables sampled from the posterior distribution. Therefore, we could compute the log determinant of the Jacobian term in $O(D)$ time as follows:
\begin{equation}
    \phi(\textbf{z}) = h^{\prime}( \textbf{w}^{\top} \textbf{z}+b) \textbf{w}
\end{equation}
\begin{equation}
    \left\vert {\rm det} \frac{\partial f_{i}}{\partial \textbf{z}_{i-1}} \right\vert = \left\vert \rm{det}(\textbf{I} + \textbf{u} \phi(z)^{\top}) \right\vert = \left\vert 1+\textbf{u}^{\top}\phi(\textbf{z}) \right\vert
\end{equation}

Finally, the network is trained by optimizing the modified ELBO based on equation~\ref{equ:lnqz}:
\begin{equation}
    \ln p(\textbf{x} \vert \textbf{c}) \geq \mathbb{E}_{q(\textbf{z}_{0} \vert \textbf{x}, \textbf{c})} \left\lbrack \ln p(\textbf{x} \vert \textbf{z}_{i}, \textbf{c}) + \sum^{i} \ln \left\vert {\rm det} \frac{\partial f_{i}}{\partial \textbf{z}_{i-1}} \right\vert \right\rbrack - {\rm KL}(q(\textbf{z}_{0} \vert \textbf{x}, \textbf{c}) \| p(\textbf{z}_i))
\end{equation}
where, $\ln p(\textbf{x} \vert \textbf{c})$ is the marginal log-likelihood of the observed data $\textbf{x}$ (i.e. here $\mathbf{x}$ represents an LV graph/mesh), conditioned on the covariates of interest (i.e. patient demographics and clinical measurements) $\textbf{c}$; $i$ is the steps of the normalizing flows. $p(\textbf{x}\vert\textbf{z}_{i}, \textbf{c})$ is the data likelihood parameterised by the decoder network, which reconstructs/predicts $\textbf{x}$ given the latent variables $\textbf{z}_{i}$, transformed by latent (planar) normalizing flows, and the conditioning variables $\textbf{c}$; ${\rm KL}(q(\textbf{z}_{0} \vert \textbf{x}) \| p(\textbf{z}_i))$ is the Kullback-Leibler divergence of the initial approximate posterior $q(\textbf{z}_{0} \vert \textbf{x}, \textbf{c})$ from the prior, $p(z)=\mathcal{N}(z \mid 0, I)$.

\section{Experimental setup and  Results}

% \subsection{Datasets}
\textbf{Data:} In this study, we created a cohort of 2360 left ventricle (LV) triangular surface meshes based on a subset of the cardiac cine-MR imaging data available from the UK Biobank (UKBB) by registering a cardiac LV atlas mesh~\cite{rodero2021linking} to the manual contours (as described in \cite{xia2022automatic}). We randomly split the dataset into 422/59/1879 for training, validation, and testing, respectively. All meshes have the same and fixed graph topology, sharing the same edges and faces but differing in the position of vertices, i.e. there is point-wise correspondence across all shapes. We used 14 covariates available for the same subjects in UKBB as conditioning variables for our model, including, gender, age, height, weight, pulse, alcohol drinker status, smoking status, HbA1c, cholesterol, C-reactive protein, glucose, high-density lipoprotein (HDL) cholesterol, insulin-like growth factor 1 (IGF-1), and low-density lipoprotein (LDL) cholesterol. These covariates were chosen as they are known cardiovascular risk factors.
% \subsection{Implementation Details}

\textbf{Implementation Details:} The framework was implemented using PyTorch on a standard PC with an NVIDIA RTX 2080Ti GPU. We trained our model using the AdamW optimizer with an initial learning rate of 1e-3 and batch size of 16 for 1000 epochs. The feature number of each graph convolutional block in the encoder was 16, 32, 32, 64, 64, and in reversed order in the decoder. The latent dimension was set as 16. The down-/up-sampling factor was four, and we used a warm-up strategy~\cite{sonderby2016ladder} to the weight of KL loss to prevent model collapse.
% \subsection{Criteria}

\textbf{Evaluation metrics:} We compared our model (cVAE-NF) with a traditional PCA-based SSM, two generative models without conditioning information including a vanilla VAE and a VAE with normalizing flow (VAE-NF), and the vanilla cVAE. We evaluated the performance of all methods using three different metrics: 1) the reconstruction error, which evaluates the generalizability of the trained model to reconstruct/represent unseen shapes, using the distance between reconstructed mesh with ground truth/original mesh; 2) the specificity error, which measures the anatomical plausibility of the synthesized virtual cohorts, using the distance between the generated meshes and its nearest neighbor in the unseen real population~\cite{davies2009building}; and 3) the variability in the left ventricular volume in the synthesized cohorts, to assess the diversity of the instances generated in terms of a clinical relevant cardiac index. LV volume variability was quantified as the standard deviation of the blood pool volumes of LV (BPVols). The Euclidean distance was used to evaluate all three metrics. Additionally, we measured the activity of the latent dimension using the statistic $A = {\rm Cov}_{\textbf{x}}(\mathbb{E}_{\textbf{z}\sim q(\textbf{z} \vert \textbf{x})} [z])$ of the observations $\textbf{x}$~\cite{burda2015importance}. A higher activity score indicates that a given latent dimension can capture greater population-wide shape variability.

% \subsection{Results}
\begin{table}[ht]
\centering
\caption{The quantitative results of the investigated methods in a hold-out test dataset. The bold values represent the results are significantly better than those of other methods.}
\label{tab:quantitative}
\begin{tabular}{c|c|c|c}
\hline
\hline
Methods &Reconstruction Error$\downarrow$ &Specificity Error$\downarrow$ &Volume Variability$\uparrow$ \\
\hline
PCA  & \textbf{0.82$\pm$0.16} &1.48$\pm$0.26  &\textbf{32.74}        \\
VAE  & 1.29$\pm$0.21 &\textbf{1.39$\pm$0.98}  &3.00         \\
VAE-NF &0.90$\pm$1.76&1.60$\pm$0.34  &16.03           \\
\hline
cVAE &1.43$\pm$0.26 &\textbf{1.32$\pm$0.21} &28.39      \\
Ours &\textbf{1.23$\pm$0.23} &1.38$\pm$0.20 &\textbf{29.91}      \\
\hline
\hline
\end{tabular}
\end{table}
\textbf{Results:} As shown in Table~\ref{tab:quantitative}, our method outperforms the cVAE in terms of reconstruction error and LV-volume variability captured in the synthesised VP (the reference volume variability for the real UKBB population used was 33.38$mm^3$). The cVAE, however, achieved lower specificity errors than our model. This indicates the superior ability of our method to capture the shape variability of the population, at the cost of generating some instances that are farther away from the real population (than those generated by cVAE), resulting in higher specificity errors. We attribute this to the ability of the normalizing flow to learn a more flexible approximate posterior latent distribution of the observed shapes, than the cVAE. This is also observed by comparing the performance of VAE and VAE-NF, where the latter can synthesize significantly more diverse VPs (e.g. it improves the volume variability from 3.00 to 16.03). Figure~\ref{fig:activity} describes the captured variability over each latent dimension. We observe that the VAE-NF obtains higher activity scores over all the latent dimensions compared to the vanilla cVAE. The normalizing flow allows the approximation of multi-modal latent distributions in the generative model, capturing greater shape variability. Although PCA outperforms our method in terms of generalization error and volume variability captured, it does not permit controllable synthesis of VPs based on relevant patient demographic information and clinical measurements, thus making it less useful for our application of synthesising VPs for use in ISTs.

\begin{figure*}[h]
	\centering
	\includegraphics[width=1\linewidth]{KDEandActivity-cropped.pdf}
 	% \includegraphics[width=0.6\linewidth]{KDE.png}	
	\caption{Left: Comparison of the activity scores over different latent dimensions between the cVAE and cVAE-NF; Right: Kernel density plots for BPVol from the VPs generated by cVAE and cVAE-NF and the real patient population (UKBB).}
	\label{fig:activity}
\end{figure*}

% \begin{figure*}[h]
% 	\centering
% 	\includegraphics[width=1\linewidth]{KDE.png}	
% 	\caption{Kernel distribution estimation plot of BPVol from the VPs generated by cVAE and cVAE-NF and the real patient population (UKBB).}
% 	\label{fig:kde}
% \end{figure*}

Capturing the distribution of clinically relevant biomarkers (e.g. BPVol) in the synthesised VPs, based on specified covariates/conditioning information available for real patients, is essential to effectively impose and emulate inclusion/exclusion criteria used during trial design in ISTs. Similarly, observed trends in the distribution of biomarkers according to patient sub-groups should be preserved in the synthesised VPs, relative to real populations e.g., the BPVol of women is known to be less than that of men~\cite{st2022sex}. To verify this, we generated VPs using cVAE and our method, conditioned on real patient data (covariates) from the UK Biobank. The virtual populations are divided into two subpopulations based on gender. Figure~\ref{fig:activity} summarises the BPVol distribution for both genders in the synthesised VPs and the real UKBB population, and the former capture the known trend of women having lower BPVol than men. Compared with cVAE, our model generates a VP which more closely matches the distribution of the LV blood pool volume observed in the real population. We also visualize the effect of manipulating individual attributes on two real patients in Fig~\ref{fig:manipulation}. We select two representative attributes that are significantly associated with BPVol and Myocardial volume (MyoVol): weight and age. We observe that the BPVol and MyoVol of the LV are positively correlated to the weight of patients (as expected). On the other hand, increasing the individual's age results in a smaller BPVol, but increased MyoVol (as visualized in Fig~\ref{fig:manipulation}), known to be due to cardiac hypertrophy caused by aging~\cite{chiao2015aging}.


\begin{figure*}[!tbp]
	\centering
	\includegraphics[width=0.9\linewidth]{VisualPatient-cropped.pdf}
	\caption{Two representative examples of the reconstructed shapes and their variations through manipulation over two demographic attributes, i.e., weight and Age. MyoVol and BPVol are shown in the bottom right corner.}
	\label{fig:manipulation}
\end{figure*}

\section{Conclusion}
We proposed the first conditional flow VAE model for \textit{controllable} synthesis of VPs of anatomy. The proposed approach was demonstrated to boost the flexibility of the learned latent distribution, yielding VPs that captured greater variability in LV shape compared to the vanilla cVAE. Additionally, our model was able to successfully model the relationship between covariates/conditioning variables and LV shape and was demonstrated to synthesise target VPs that fit desired criteria (in terms of patient demographics and clinical measurements) and closely matched the real population in terms of a clinically relevant biomarker (LV BPVol). These results indicate that our approach shows promise for \textit{controllable} synthesis of diverse, yet plausible VPs of anatomy. Future work will look to expand to modelling the whole-heart and will investigate the impact of individual covariates on VP synthesis in more depth. 

\section{Acknowledgement}
This research was conducted using data from the UK Biobank (access application 11350). This work was supported by the Royal Academy of Engineering (INSILEX CiET1819/19), Engineering and Physical Sciences Research Council (EPSRC) programs TUSCA EP/V04799X/1, and the Royal Society Exchange Programme CROSSLINK  IES$\backslash$NSFC$\backslash$201380.
%
% ---- Bibliography ----
%
% BibTeX users should specify bibliography style 'splncs04'.
% References will then be sorted and formatted in the correct style.
%
\bibliographystyle{splncs04}
\bibliography{ref}
%

\clearpage
\includepdfmerge{SupplyMaterial.pdf, 1}
\includepdfmerge{SupplyMaterial.pdf, 2}

\end{document}
